\section{Verificação experimental}
\label{sec:verificacao}

Para confirmar o entendimento descrito nas seções anteriores, escrevemos um programa de teste
(\code{tools/driver.cpp}, \cref{app:repro}) que usa os arquivos \code{IBHM.cpp} e \code{IBHMInterfaceElementBuilder.cpp}
como estão, insere arestas (uma a uma ou em lote), registra o operador escolhido em cada extremidade e, após cada
inserção, verifica:
\begin{enumerate}[label=\textbf{V\arabic*.}]
\item simetria das opostas: $\Op[\Op[h]]=h$ para toda oposta válida;
\item consistência de vértices entre opostas: $\V[h]=\V[\nxt(\Op[h])]$ e $\V[\nxt(h)]=\V[\Op[h]]$;
\item toda half-edge \COLLAPSED{} tem comprimento zero: $\V[h]=\V[\nxt(h)]$;
\item variedade: os cantos de cada vértice formam um único leque conexo pela relação $h\sim\Op[\prv(h)]$ (nenhum vértice
``estrangulado'' por dois leques separados);
\item todo vértice criado é usado por algum elemento;
\item o número de vértices criados é igual ao previsto por \code{computeNumberOfNewVertices};
\item toda cópia tem exatamente as coordenadas do seu vértice de origem.
\end{enumerate}
As verificações V1--V5 e V7 ficam em \code{tests/support/MeshValidation} e são as mesmas usadas pelos testes
automatizados (\cref{sec:testes}).
Para comparar resultados de ordens diferentes, usamos a partição dos cantos dos elementos \emph{originais} em classes de
vértices: duas execuções produzem a mesma topologia (a menos de renumeração) se e somente se essas partições coincidem.

\subsection{Malhas aleatórias}

O teste com malhas aleatórias (\code{tools/stress.py}) gera grades de $3$ a $7$ por $3$ a $7$ células de quatro tipos
--- só quadriláteros, só triângulos, quadriláteros e triângulos misturados, e misturadas com hexágonos (dois
quadriláteros fundidos) ---, sorteia até metade das arestas internas e insere esse conjunto quatro vezes: em três ordens
aleatórias, escolhendo ao acaso qual half-edge representa cada aresta, e uma vez em lote, numa única chamada. As quatro
execuções precisam ser válidas e produzir a mesma topologia. Foram feitas duas rodadas de 400 malhas
(\cref{tab:estresse}), refeitas com o código desta versão: os números são os mesmos da v0.3, agora também com a V7.

\begin{table}[ht]
\centering\small
\begin{tabularx}{\linewidth}{>{\raggedright}X>{\centering}p{3.3cm}>{\centering\arraybackslash}p{3.3cm}}
\toprule
& Rodada 1: sem restrição & Rodada 2: sem vértice de borda com 2+ arestas\\
\midrule
Malhas (quad / tri / híbrida / polígonos) & 100 / 100 / 100 / 100 & 100 / 100 / 100 / 100\\
Concluídas e válidas (V1--V7) & 259 & \textbf{400}\\
\quad por tipo & 100 / 37 / 56 / 66 & 100 / 100 / 100 / 100\\
Abortadas & 141 & 0\\
\quad das quais com vértice de borda com 2+ arestas & \textbf{141} & ---\\
\quad asserção em \code{openHole} / em \code{heNext} & 72 / 69 & ---\\
Ordens diferentes ou lote com resultado diferente & 0 & 0\\
Contagem de vértices diferente da prevista & 0 & 0\\
Operadores aplicados (Canônico / Split / Expand / Hole) & 3124 / 553 / 384 / 321 & 5130 / 791 / 708 / 535\\
\bottomrule
\end{tabularx}
\caption{Teste com malhas aleatórias. Todas as falhas da rodada 1 são exatamente os casos em que algum vértice de borda
recebe duas ou mais arestas. Nas malhas só de quadriláteros essa configuração não aparece (em uma grade, cada vértice de
borda tem uma única aresta interna); nas de triângulos ela é frequente. As asserções que interrompem a execução são a
do primeiro laço do \code{openHole} e a de \code{heNext}, chamada pelo segundo laço com \BORDER{}
(\cref{sec:falha-borda}).}
\label{tab:estresse}
\end{table}

\subsection{Casos específicos}

A \cref{tab:verificacao} reúne os cenários construídos à mão: o exemplo dos slides em várias ordens, uma trinca reta e
um laço fechado, o experimento que troca Split Element por Expand Edge e um conjunto de entradas inválidas. As linhas
marcadas com \textdagger{} mudaram com as correções desde a v0.3 (I3, I4 e I5).

\begin{table}[ht]
\centering\small
\begin{tabularx}{\linewidth}{>{\raggedright}p{5.6cm}X}
\toprule
Cenário & Resultado\\
\midrule
Exemplo dos slides (ordem dos slides) & válido em todos os passos; 24/24 operadores iguais aos slides; 19 vértices novos
(previsto: 19); buracos em $K$, $L$ e $M$.\\
Exemplo, ordem do \code{main.cpp}, lote único e mais 21 ordens & mesma topologia da ordem dos slides.\\
Coordenadas das cópias (V7)\,\textdagger & toda cópia tem as coordenadas do original, no exemplo e nos 659 casos
aleatórios válidos (até a v0.3 as cópias nasciam em $(0,0)$).\\
Trinca reta $J$--$K$--$L$--$M$--$N$ em 3 ordens & válida; mesma topologia nas 3 ordens.\\
Laço fechado em torno do elemento $FGLK$ & válido; Split Element nos quatro cantos.\\
Split Element substituído por Expand Edge\textsuperscript{a} & 791 ocorrências nas 400 malhas da rodada 2: topologia
idêntica nas 1600 execuções.\\
Mesma aresta duas vezes (em chamadas separadas ou no mesmo lote)\,\textdagger & a segunda é rejeitada; a pré-contagem é
exata (antes reservava vértices que não eram usados).\\
Aresta de borda na entrada\,\textdagger & rejeitada; contada como 0 vértices (antes, 2).\\
Half-edge de um elemento de interface na entrada\,\textdagger & rejeitada, sem alterar a malha (antes, \textbf{laço
infinito} em \code{isBorder}, dentro de \code{computeNumberOfNewVertices}).\\
Índice de half-edge maior que o vetor\,\textdagger & rejeitado, sem alterar a malha (antes, asserção em
\code{heOpposite}; em \emph{release}, leitura fora do vetor).\\
Índice no espaço reservado e ainda não validado\textsuperscript{b}\,\textdagger & rejeitado; a posição reservada de
$\Op$ vale \BORDER{} (antes valia 0, a aresta era aceita e a inserção terminava numa asserção do Open Hole).\\
Malha com furo interno (\code{meshIBHM.txt}): aresta $0$--$4$ e depois $4$--$8$ & a primeira é válida (Canônico de
borda nas duas pontas); a segunda falha no Open Hole, porque o vértice 4 está na borda do furo (I1).\\
\bottomrule
\end{tabularx}
\caption{Cenários específicos. \textsuperscript{a}Variante compilada à parte em que \code{retrieveOperator} nunca
escolhe Split Element. \textsuperscript{b}Por exemplo, a half-edge 48 de \code{malha2Hyb.txt} depois de uma chamada
cuja única aresta foi rejeitada: a reserva criou as posições 48--51, mas nenhum elemento foi validado.}
\label{tab:verificacao}
\end{table}

Em resumo: \emph{para todas as configurações suportadas, a implementação produz topologias válidas, independentes da
ordem e do modo de inserção, com contagem exata de vértices e cópias com as coordenadas do original}. Desde a v0.3, as
entradas inválidas passaram a ser rejeitadas sem alterar a malha. A falha que resta é o caso de borda ainda não tratado,
discutido no fim desta seção e na \cref{sec:defeitos}.

\subsection{Testes automatizados}
\label{sec:testes}

Os cenários acima viraram testes automatizados do repositório (\code{tests/}), escritos com o Google Test e executados
pelo CTest (\cref{app:repro}). Cada teste roda em um processo próprio, com limite de tempo, de modo que uma asserção da
biblioteca ou um laço infinito faz falhar só aquele teste. O núcleo é compilado de novo para os testes com as asserções
ligadas, qualquer que seja o tipo de build, porque alguns defeitos só aparecem como asserção. A \cref{tab:suites} lista
as suítes; o teste com malhas aleatórias também é registrado no CTest.

\begin{table}[ht]
\centering\small
\begin{tabularx}{\linewidth}{>{\ttfamily\small}p{3.6cm}c X}
\toprule
\normalfont Suíte & Testes & O que verifica\\
\midrule
test\_ibhm & 13 & reserva exata e inclusão de pontos, \code{duplicatePoint} (inclusive com realocação durante a cópia),
3 coordenadas, um teste de morte (\code{addPoint} com número errado de coordenadas dispara a asserção),
\code{numberOfHalfEdges} e $\Op[\,]$ reservado igual a \BORDER.\\
test\_vertex\_count & 6 & a pré-contagem: exemplo em lote, aresta repetida, aresta de borda, vértice já em elemento de
interface, vértice de borda com duas arestas (só a contagem).\\
test\_base\_mesh & 5 & o exemplo dos slides: os 24 operadores e os vértices novos de cada inserção, com V1--V7 a cada
passo; malha final; lote igual a uma por vez; duas chamadas iguais a uma; 21 ordens com a half-edge sorteada.\\
test\_invalid\_input & 7 & rejeições sem alterar a malha: índice fora do vetor, índice no espaço reservado, índice de
elemento criado na mesma chamada, aresta de borda, aresta de elemento de interface, aresta já aberta, aresta repetida.\\
test\_known\_defects & 3 & desativados: I1 (uma aresta por chamada e em lote) e I2 (\code{isBorder}). Descrevem o
comportamento correto; hoje, rodados à força, o I1 aborta no Open Hole e o I2 não termina.\\
stress\_supported & 1 & \code{stress.py} com a semente 22 e 400 malhas, sem a configuração do I1 (rodada 2).\\
stress\_all & 1 & rodada 1 (semente 21, todas as configurações); desativado até o I1 ser corrigido.\\
\bottomrule
\end{tabularx}
\caption{Testes automatizados. Os 31 testes ativos das suítes do Google Test levam cerca de 0,2\,s; o
\code{stress\_supported}, cerca de 7\,s.}
\label{tab:suites}
\end{table}

Para saber se os testes detectam os defeitos que já foram corrigidos, as mesmas suítes foram compiladas contra o núcleo
de versões anteriores (\cref{tab:controle}). Cada correção é detectada por mais de um teste, e o núcleo atual passa em
todos.

\begin{table}[ht]
\centering\small
\begin{tabularx}{\linewidth}{>{\raggedright}p{4.6cm}c X}
\toprule
Núcleo & Falhas & Testes que falham\\
\midrule
atual & 0 de 31 & ---\\
antes do I3 (commit \code{3b991e3}) & 4 de 31 & $\Op[\,]$ reservado; índice fora do vetor e índice de elemento da mesma
chamada (asserção em \code{heOpposite}); índice no espaço reservado (asserção no Open Hole).\\
antes do I5 (\code{2a4eb0e}) & 12 de 31 & os anteriores, cinco testes de contagem e três de entrada inválida, sendo
três por laço infinito.\\
antes do I4 (\code{28dd85d}) & 15 de 18\textsuperscript{*} & os mesmos testes de contagem e de entrada inválida e os
quatro testes do exemplo que verificam a V7.\\
\bottomrule
\end{tabularx}
\caption{Controle negativo: as suítes compiladas contra o núcleo de versões anteriores. \textsuperscript{*}Sem
\code{test\_ibhm}, que usa \code{duplicatePoint}, inexistente nessa versão.}
\label{tab:controle}
\end{table}

\subsection{O caso não suportado: vértice de borda com duas arestas}
\label{sec:falha-borda}

Quando um vértice de borda já foi aberto pelo Canônico e recebe uma segunda aresta, ele pertence a um elemento de
interface e não tem aresta colapsada; \code{retrieveOperator} escolhe então Open Hole. Mas um dos dois laços de
\code{openHole} chega à borda da malha antes de encontrar um elemento de interface (\cref{fig:borda}). Em modo
\emph{debug}, a asserção \code{currentHE != IBHM::BORDER} interrompe o programa (os comentários do código perguntam
``\emph{@todo Is it possible?}'': sim, é possível). Em modo \emph{release} o resultado é pior: o primeiro laço escreve
na posição \code{0xFFFFFFFF} do vetor de vértices, e o segundo laço não detecta a borda, porque testa $\nxt(\Op[h])$ em
vez de $\Op[h]$ --- e, na IBHM, $\nxt(\BORDER)$ é uma busca binária com um índice inválido, cujo resultado é
imprevisível.

Esse caso não é raro: aparece sempre que uma trinca chega à borda por um vértice que já tem outra aresta selecionada, o
que é comum em malhas de triângulos (63 das 100 malhas de triângulos da rodada 1) e nos vértices da borda de furos
internos.

\begin{figure}[ht]
\centering
\begin{tikzpicture}[scale=2.200, every node/.style={font=\scriptsize}]
\filldraw[cefill] (0.130,0.075) -- (1.000,0.000) -- (1.200,0.700) -- (0.552,0.725) -- cycle;
\filldraw[cefill] (-0.075,0.130) -- (0.352,0.841) -- (0.000,1.400) -- (-0.500,0.866) -- cycle;
\filldraw[cefill] (-0.075,0.130) -- (-0.500,0.866) -- (-1.200,0.700) -- (-1.000,0.000) -- cycle;
\filldraw[iefill] (0.130,0.075) -- (0.552,0.725) -- (0.352,0.841) -- (-0.075,0.130) -- cycle;
\draw[seledge] (-0.075,0.130) -- (-0.500,0.866);
\draw[hearrow, red!70!black] (-0.306,0.267) -- (-0.544,0.679);
\node[helabel, text=red!70!black] at (-0.549,0.450) {$a$};
\node[vie, minimum size=0.29cm] at (0.130,0.075) {$Q$};
\node[vreg, dot] at (1.000,0.000) {};
\node[vie, minimum size=0.29cm] at (0.352,0.841) {$R$};
\node[vreg, minimum size=0.29cm] at (-0.500,0.866) {$P$};
\node[vreg, dot] at (-1.000,0.000) {};
\node[vreg, dot] at (1.200,0.700) {};
\node[vreg, dot] at (0.000,1.400) {};
\node[vreg, dot] at (-1.200,0.700) {};
\node[vie, minimum size=0.29cm] at (-0.075,0.130) {$Q_{1}$};
\node[vie, minimum size=0.29cm] at (0.552,0.725) {$R_{1}$};

\end{tikzpicture}

\caption{Configuração que falha. $Q$ está na borda e já foi aberto pela aresta $QR$ (cópias $Q$ e $Q_1$; os dois
elementos à esquerda de $QR$ passaram a usar $Q_1$). Ao inserir $QP$, o Open Hole em $Q_1$ gira no sentido
anti-horário a partir de $he$ e alcança imediatamente a borda (à esquerda), sem encontrar elemento de interface. Nesse
caso não há buraco a abrir: bastaria duplicar $Q_1$ e mover para a cópia nova os elementos entre $QP$ e a borda (como no
Canônico de borda), deixando o lado curto do novo elemento na borda.}
\label{fig:borda}
\end{figure}
